%% ============================================================
%%  Chapter 3 — Methodology
%% ============================================================
\chapter{Methodology}
\label{ch:methodology}

\section{System Overview}

HoneyIQ is a closed-loop simulation in which an \emph{Attacker} and a
\emph{Defender} interact through a Gymnasium environment
(\texttt{CyberSecurityEnv}).
Figure~\ref{fig:architecture} shows the component layout.

\begin{figure}[ht]
  \centering
  % Replace with \includegraphics when a diagram is available
  \framebox[0.85\textwidth]{%
    \parbox{0.82\textwidth}{\centering\vspace{2cm}
    [Architecture diagram placeholder — see docs/architecture.md]
    \vspace{2cm}}%
  }
  \caption{High-level architecture of HoneyIQ.}
  \label{fig:architecture}
\end{figure}

At each time step the attacker generates an attack type, kill-chain stage,
and feature vector; the environment computes a threat level and state vector;
the defender selects a honeypot action; and a reward is returned to the
defender's replay buffer.

\section{Attacker Module}

\subsection{Attack Type and Kill Chain Enumerations}

Two \texttt{IntEnum} classes define the discrete state spaces:

\begin{itemize}
  \item \texttt{AttackType} --- ten categories (NORMAL through WORMS) with
        associated severity weights in $[0, 0.90]$.
  \item \texttt{KillChainStage} --- seven stages mapped to the Lockheed
        Martin Kill Chain.
\end{itemize}

Table~\ref{tab:attack_severity} lists each attack type, its severity weight,
and its typical kill-chain stage.

\begin{table}[ht]
  \centering
  \caption{Attack types, severity weights, and primary kill-chain stages.}
  \label{tab:attack_severity}
  \begin{tabular}{clcc}
    \toprule
    Index & Attack Type & Severity & Primary Stage \\
    \midrule
    0 & NORMAL        & 0.00 & Reconnaissance \\
    1 & RECONNAISSANCE & 0.20 & Reconnaissance \\
    2 & ANALYSIS      & 0.25 & Weaponization \\
    3 & FUZZERS       & 0.35 & Delivery \\
    4 & GENERIC       & 0.40 & Delivery \\
    5 & EXPLOITS      & 0.70 & Exploitation \\
    6 & SHELLCODE     & 0.75 & Exploitation \\
    7 & BACKDOORS     & 0.80 & Installation \\
    8 & DOS           & 0.85 & Actions on Objectives \\
    9 & WORMS         & 0.90 & Command \& Control \\
    \bottomrule
  \end{tabular}
\end{table}

\subsection{Markov Transition Model}

The attacker maintains two Discrete-Time Markov Chains.
The $10 \times 10$ attack-type matrix $\mathbf{T}_A$ and the $7 \times 7$
kill-chain-stage matrix $\mathbf{T}_K$ are both row-stochastic.
At each step the attacker samples:
%
\begin{align}
  a_{t+1} &\sim \text{Categorical}\bigl(\mathbf{T}_A[a_t, \cdot]\bigr), \\
  k_{t+1} &\sim \text{Categorical}\bigl(\mathbf{T}_K[k_t, \cdot]\bigr).
\end{align}
%
A kill-chain constraint prevents unrealistic regression by flooring the
sampled next stage:
%
\begin{equation}
  k_{t+1} \leftarrow \max\bigl(\tilde{k}_{t+1},\;
                                 \text{primary\_stage}(a_{t+1}) - 1\bigr).
\end{equation}

\subsubsection{Intent Profiles}

Four intent modifiers skew the base matrices (Table~\ref{tab:intents}).

\begin{table}[ht]
  \centering
  \caption{Attacker intent profiles and their behavioural characteristics.}
  \label{tab:intents}
  \begin{tabular}{lll}
    \toprule
    Intent & Attack preference & Kill chain speed \\
    \midrule
    STEALTHY      & Recon, Analysis, Backdoors & Slow; high self-loop probability \\
    AGGRESSIVE    & DoS, Worms, Exploits       & Fast; low self-loops, forward bias \\
    TARGETED      & Exploits $\to$ Shellcode $\to$ Backdoors & Direct path through exploitation \\
    OPPORTUNISTIC & Generic, Fuzzers (scattered) & Moderate forward bias \\
    \bottomrule
  \end{tabular}
\end{table}

\subsection{Network Feature Simulation}

At each step the attacker samples 15 network features from parametric
distributions conditioned on the current attack type.
Selected signatures are:
%
\begin{itemize}
  \item \textbf{Reconnaissance}: very short \texttt{dur}, high
        \texttt{ct\_dst\_ltm} (port-scanning pattern).
  \item \textbf{DoS}: extreme \texttt{sload} (50k--500k bps), near-zero
        \texttt{dur}, Poisson-distributed \texttt{spkts} ($\lambda=500$).
  \item \textbf{Backdoors}: long \texttt{dur} (10--3600\,s), low
        \texttt{sload} (stealthy), balanced window sizes.
  \item \textbf{Worms}: high \texttt{ct\_dst\_ltm} ($\lambda=40$), high
        packet counts (spreading behaviour).
\end{itemize}

\section{Defender Module}

\subsection{Attack Classifier}

A Random Forest classifier is trained on 600 synthetic samples per class
(6\,000 total) generated from the feature distributions described above.
Features are standardised with \texttt{StandardScaler}.
The classifier provides a predicted attack type at each step for logging and
introspection; in the current implementation the DQN state uses ground-truth
labels, but the classifier output would replace these in a real deployment.

Classifier performance is evaluated on a held-out set (seed = 999) generated
independently of the training split, reporting per-class precision, recall,
F1-score, and overall accuracy.

\subsection{DQN Agent}

The DQN agent (\texttt{DQNAgent}) wraps the policy network, target network,
replay buffer, and training loop.
Table~\ref{tab:hyperparams} summarises the key hyperparameters.

\begin{table}[ht]
  \centering
  \caption{DQN hyperparameters.}
  \label{tab:hyperparams}
  \begin{tabular}{ll}
    \toprule
    Parameter & Value \\
    \midrule
    State dimension & 24 \\
    Action space & 5 (HoneypotAction) \\
    Hidden layers & [256, 128, 64] \\
    Normalisation & LayerNorm \\
    Activation & ReLU \\
    Replay buffer capacity & 15\,000 \\
    Mini-batch size & 64 \\
    Discount factor $\gamma$ & 0.99 \\
    Learning rate & $10^{-3}$ \\
    Target network update interval & 150 steps \\
    $\varepsilon$ initial / final & 1.0 / 0.05 \\
    $\varepsilon$ decay & 0.997 \\
    Gradient clip norm & 10.0 \\
    Loss function & Huber (SmoothL1) \\
    \bottomrule
  \end{tabular}
\end{table}

\subsection{Honeypot Action Space and Reward Function}

The reward function $R(s, a)$ encodes domain knowledge as a base matrix
indexed by action and threat band, with three layers of modifiers:

\begin{enumerate}
  \item \textbf{Base matrix} (Table~\ref{tab:reward_matrix}) --- values
        range from $-6$ (ALLOW on critical traffic) to $+6$ (ALERT on
        critical traffic).
  \item \textbf{Late kill-chain penalty} --- negative rewards are multiplied
        by $1.5\times$ for stages $\geq$ INSTALLATION.
  \item \textbf{Attack-type bonuses} --- TROLL on BACKDOORS/SHELLCODE/WORMS
        ($+0.8$); BLOCK on WORMS ($+1.0$); LOG on RECONNAISSANCE ($+0.5$).
\end{enumerate}

\begin{table}[ht]
  \centering
  \caption{Base reward matrix (action $\times$ threat band).}
  \label{tab:reward_matrix}
  \begin{tabular}{lccccc}
    \toprule
    Action & Benign & Low & Medium & High & Critical \\
    \midrule
    ALLOW & $+1.0$ & $-1.0$ & $-2.0$ & $-4.0$ & $-6.0$ \\
    LOG   & $-0.5$ & $+1.5$ & $+2.0$ & $+0.5$ & $-1.0$ \\
    TROLL & $-1.0$ & $+0.5$ & $+1.0$ & $+3.0$ & $+2.0$ \\
    BLOCK & $-2.0$ & $-0.5$ & $+0.5$ & $+4.0$ & $+5.0$ \\
    ALERT & $-3.0$ & $-1.0$ & $-0.5$ & $+2.0$ & $+6.0$ \\
    \bottomrule
  \end{tabular}
\end{table}

\section{Environment: \texttt{CyberSecurityEnv}}

\texttt{CyberSecurityEnv} extends \texttt{gymnasium.Env} and bridges the
attacker and defender.
It exposes:

\begin{itemize}
  \item \textbf{Observation space}: \texttt{Box}(24,) with features
        described in Section~\ref{sec:state_encoding}.
  \item \textbf{Action space}: \texttt{Discrete}(5).
  \item \textbf{Episode lifecycle}: \texttt{reset()} re-initialises the
        attacker and clears escalation history; each \texttt{step()} call
        advances the attacker, constructs the state vector, computes the
        threat level and reward, and increments the step counter.
        Episodes truncate at a configurable \texttt{max\_steps}.
\end{itemize}

\subsection{State Vector Encoding}
\label{sec:state_encoding}

The 24-dimensional observation vector uses one-hot encoding for categorical
variables and normalised scalars for continuous ones:

\begin{equation}
  s = \underbrace{[\,\text{attack\_type}_{1:10}\,]}_{\text{one-hot}}
      \;\oplus\;
      \underbrace{[\,\text{stage}_{1:7}\,]}_{\text{one-hot}}
      \;\oplus\;
      \underbrace{[\,\text{threat},\;\tfrac{\text{count}}{100},\;\text{esc\_rate}\,]}_{\text{scaled scalars}}
      \;\oplus\;
      \underbrace{[\,\text{intent}_{1:4}\,]}_{\text{one-hot}}.
\end{equation}

\subsection{Composite Threat Level}
\label{sec:threat_level}

The threat level $T \in [0, 1]$ combines four signals:

\begin{equation}
  T = 0.45 \cdot s_{\text{attack}}
    + 0.35 \cdot w_{\text{stage}}
    + 0.15 \cdot r_{\text{esc}}
    + 0.05 \cdot \min\!\left(1,\, \frac{n_{\text{attack}}}{100}\right),
\end{equation}

where $s_{\text{attack}}$ is the per-type severity, $w_{\text{stage}}$ is
the kill-chain stage weight, $r_{\text{esc}}$ is the fraction of attacks in
the last 20 steps, and $n_{\text{attack}}$ is the cumulative attack count.
Threat bands partition $[0,1]$ into five tiers: Benign ($<0.15$), Low
(0.15--0.35), Medium (0.35--0.55), High (0.55--0.75), and Critical
($\geq 0.75$).

\section{Evaluation Infrastructure}

\subsection{Metrics}

Two granularity levels are maintained:

\begin{itemize}
  \item \textbf{\texttt{StepRecord}} --- per-step tuple:
        \texttt{(episode, step, action, reward, attack\_type,
        kill\_chain\_stage, threat\_level, is\_attack,
        predicted\_attack, loss, escalation\_rate)}.
  \item \textbf{\texttt{EpisodeRecord}} --- per-episode aggregates:
        total reward, detection rate, false-positive rate, average threat
        level, average DQN loss, action distribution, kill-chain
        distribution.
\end{itemize}

The detection confusion matrix is defined relative to the ALLOW action:
any action other than ALLOW on an attack step counts as a true positive.

\begin{align}
  \text{detection\_rate}      &= \frac{\text{TP}}{\text{TP} + \text{FN}}, \\
  \text{false\_positive\_rate} &= \frac{\text{FP}}{\text{FP} + \text{TN}}.
\end{align}

\subsection{Experimental Protocol}

All experiments use the same random seed for reproducibility.
Training runs for a configurable number of episodes (default: 500) using
the \emph{Opportunistic} intent profile.
Cross-intent evaluation freezes the trained policy and re-evaluates it
against each of the four intent profiles for 50 episodes each, recording
mean detection rate and false-positive rate.
